% conclusion：生成于阅读卡+证据；引用检查通过


\section{结论}

现有长期记忆系统的基础机制可归纳为三条技术脉络。第一条是以生成式智能体架构为代表的“记忆流+检索+反思”范式，通过自然语言记录经验，并在检索时综合相关性、新近性与重要性，再以反思模块生成高层推断\cite{W0001}。第二条是虚拟上下文管理路线，以 MemGPT 为代表，将主上下文划分为只读系统指令、可写工作上下文与 FIFO 消息队列，并借助函数调用实现对外部上下文的自主读写与驱逐\cite{W0002}。第三条是读写分离与经验提取路线：RET-LLM 通过微调生成记忆 API 调用，在不修改 LLM 架构的前提下实现显式存储与检索\cite{W0680}；ExpeL 通过自主收集经验并提取跨任务洞察，将回忆对象从原始经验扩展为自然语言知识\cite{W0618}。此外，MemoryBank 引入艾宾浩斯遗忘曲线启发的记忆更新机制，以记忆强度衰减模拟人类记忆的动态性\cite{W0003}；Reflexion 则以语言化自我反思构成情景经验记忆，作者报告其在 HumanEval 编码基准上以自反策略达到 91\% pass@1\cite{W0004}。综合来看，基础系统层已形成“写入—存储—检索—更新”的通用骨架，但在写入时机、记忆粒度和更新策略上仍缺乏统一约定\cite{W0001,W0002,W0680,W0618}。

\subsection{RQ1：长期记忆与相关概念的边界是什么？}

本文的操作定义将长期记忆界定为：由智能体显式写入并跨会话持久化的信息状态。这一边界排除了三类相关但不同的机制：上下文扩展仅提供会话内召回，不构成跨会话的持久写入\cite{W0039}；静态检索增强的外部知识库若无智能体写入行为且不随交互演化，不进入记忆生命周期\cite{W0039}；参数化知识在推理时无法被精确更新，也难以携带个体交互轨迹\cite{W0002,W0618}。据此，综述覆盖的工作以"写入—存储—检索—使用—更新"的显式闭环为共同特征，而机制形态（记忆流、图结构、分层摘要、经验规则）与记忆类型（事实、事件、任务状态、程序性经验）构成设计空间的两条主轴\cite{W0140,W0039}。

\subsection{RQ2：组织与检索机制如何影响记忆系统的性能？}

组织与检索机制构成了长期记忆系统性能分化的关键因素。时间知识图谱类系统（如 Zep）通过情节子图、语义实体子图与社区子图三层结构组织信息，并引入图感知重排器以提升检索质量\cite{W0684}。Mem++ 采用非破坏性组织策略，在写入时完整保存带日期与作者的文档，读取时按问题时间上界筛选，并在冲突场景中保留版本而非覆盖\cite{W0028}。分层图记忆系统 HiGMem 则通过 LLM 引导的证据集裁剪，在 LoCoMo10 基准上将平均检索轮数从 A-Mem 的 99.84 降至 8.09，同时取得更高精确率与相近召回率\cite{W0030}。扩散激活机制（Synapse）通过融合情景与语义记忆，在 LoCoMo 上取得加权平均 F1 40.5，超过 Zep 的 39.7 与 AriGraph 的 33.7\cite{W0099}。然而，组织复杂性的增益并非线性：Mem0 结合图记忆的整体得分仅比基础版高约 2\%\cite{W0682}；一项图结构与扁平向量基线的对比研究显示，图结构记忆的 token F1 为 0.417，低于扁平向量基线的 0.468\cite{W0034}。这表明，在特定评测条件下，组织机制的优势高度依赖于任务与检索设置，不能简单假定结构化组织必然优于扁平存储\cite{W0034,W0682}。

\subsection{RQ3：更新与遗忘机制如何影响长期一致性？}

更新与遗忘机制对长期一致性具有双重作用。在正向更新方面，检索驱动再巩固框架 REALM 将插入、合并与遗忘统一为单一整合阶段，在 LoCoMo 上平均准确率达 75.97\%，比最强基线高 7.17 个百分点\cite{W0012}。SodaMem 通过 supersession 机制在新事实写入时主动失效旧事实，并以软时间窗口而非硬过滤处理时序查询\cite{W0222}。在遗忘治理方面，Memory Worth 估计在 10,000 回合后与真实效用的斯皮尔曼秩相关达 ρ=0.89±0.02，显示在线价值估计具备可行性\cite{W0199}。然而，遗忘也可能削弱系统能力：选择性遗忘一次应用移除 9.8\% 节点后，judged correctness 下降 1.6 点\cite{W0034}；记忆规范化研究显示，全规范化记忆的情感一致性改善在多重比较校正后不再显著，事实漂移在任何校正水平下均不显著，作者明确将其定位为探索性发现\cite{W0077}。这表明更新与遗忘机制在提升一致性的同时，可能引入信息损失风险，需在价值估计与保留策略之间谨慎权衡\cite{W0034,W0077,W0199}。

\subsection{RQ5：现有评测框架能否揭示长期记忆的真实能力？}

现有评测框架已在覆盖维度上取得显著进展，但对长期记忆真实能力的刻画仍不充分。LoCoMo 将对话规模扩展至约 16K tokens、588.2 turns、27.2 sessions，揭示了长上下文 LLM 与 RAG 在时间与因果动态理解上的显著落后\cite{W0005}。LongMemEval 明确将知识更新与弃答纳入五大核心能力\cite{W0006}。Memora 则进一步指出，LoCoMo 中 94\% 的评估问题只需不超过 2 个此前会话的证据，LongMemEval 中该比例为 85\%，显示既有基准偏重浅层检索而缺乏记忆整合压力\cite{W0171}。成本感知评测表明，长期记忆可将依赖任务成功率从无记忆的 0.00 提升至 0.55–1.00，但嵌入检索在事实变化时表现不稳定（0.30–0.95），且代理仅在 55\% 的时间对正确检索值采取行动\cite{W0031}。操作级幻觉评测 HaluMem 则揭示，现有系统倾向于在抽取和更新阶段产生并累积幻觉，再向问答阶段传播\cite{W0261}。这些证据共同指向一个结论：当前评测正向更细粒度的记忆操作和可靠性维度演进，但“记忆何时真正帮助、何时引入危害”仍缺乏系统化的成本—收益度量\cite{W0031,W0171,W0261}。

\subsection{RQ4：可靠性与安全风险如何约束记忆系统的设计？}

可靠性与安全风险已成为长期记忆系统不可回避的设计约束。在非恶意可靠性方面，来源感知系统 Agent Zero Memory 通过引用锁纪律在 LongMemEval 上达 95.60\%，较此前最强系统提升 +0.73 分\cite{W0010}。多智能体错误传播研究显示，幻觉检测率从 Stage 1 的 72.0\% 单调下降至 Stage 4 的 50.9\%，末端检查仅改善 2.3 个百分点\cite{W0688}。信息压缩的统一速率—失真视角进一步指出，查询无关压缩的有效每查询预算会被查询分布熵削减，导致快照式方法在查询分布变化时性能塌陷\cite{W0207}。在安全投毒方面，MemPoison 显示选择性记忆机制可被绕过，攻击成功率最高达 0.95\cite{W0319}；而记忆投毒防御在良性流量中亦有代价——读取端防御使核心准确率下降 4.4 点，33.6\% 的合法记忆被误隔离\cite{W0312}。安全综述进一步将长期记忆界定为“可写、跨会话持久”的新安全基底，并指出保留策略可能意外保留对抗性条目而丢弃合法条目\cite{W0039}。综合而言，可靠性与安全约束要求记忆系统在来源验证、写入控制与良性效用之间寻找平衡，单一环节的强化不足以应对传播链上的综合风险\cite{W0010,W0312,W0039,W0688}。

